% Options for packages loaded elsewhere
% Options for packages loaded elsewhere
\PassOptionsToPackage{unicode}{hyperref}
\PassOptionsToPackage{hyphens}{url}
\PassOptionsToPackage{dvipsnames,svgnames,x11names}{xcolor}
%
\documentclass[
  letterpaper,
  DIV=11,
  numbers=noendperiod]{scrartcl}
\usepackage{xcolor}
\usepackage{amsmath,amssymb}
\setcounter{secnumdepth}{-\maxdimen} % remove section numbering
\usepackage{iftex}
\ifPDFTeX
  \usepackage[T1]{fontenc}
  \usepackage[utf8]{inputenc}
  \usepackage{textcomp} % provide euro and other symbols
\else % if luatex or xetex
  \usepackage{unicode-math} % this also loads fontspec
  \defaultfontfeatures{Scale=MatchLowercase}
  \defaultfontfeatures[\rmfamily]{Ligatures=TeX,Scale=1}
\fi
\usepackage{lmodern}
\ifPDFTeX\else
  % xetex/luatex font selection
\fi
% Use upquote if available, for straight quotes in verbatim environments
\IfFileExists{upquote.sty}{\usepackage{upquote}}{}
\IfFileExists{microtype.sty}{% use microtype if available
  \usepackage[]{microtype}
  \UseMicrotypeSet[protrusion]{basicmath} % disable protrusion for tt fonts
}{}
\makeatletter
\@ifundefined{KOMAClassName}{% if non-KOMA class
  \IfFileExists{parskip.sty}{%
    \usepackage{parskip}
  }{% else
    \setlength{\parindent}{0pt}
    \setlength{\parskip}{6pt plus 2pt minus 1pt}}
}{% if KOMA class
  \KOMAoptions{parskip=half}}
\makeatother
% Make \paragraph and \subparagraph free-standing
\makeatletter
\ifx\paragraph\undefined\else
  \let\oldparagraph\paragraph
  \renewcommand{\paragraph}{
    \@ifstar
      \xxxParagraphStar
      \xxxParagraphNoStar
  }
  \newcommand{\xxxParagraphStar}[1]{\oldparagraph*{#1}\mbox{}}
  \newcommand{\xxxParagraphNoStar}[1]{\oldparagraph{#1}\mbox{}}
\fi
\ifx\subparagraph\undefined\else
  \let\oldsubparagraph\subparagraph
  \renewcommand{\subparagraph}{
    \@ifstar
      \xxxSubParagraphStar
      \xxxSubParagraphNoStar
  }
  \newcommand{\xxxSubParagraphStar}[1]{\oldsubparagraph*{#1}\mbox{}}
  \newcommand{\xxxSubParagraphNoStar}[1]{\oldsubparagraph{#1}\mbox{}}
\fi
\makeatother


\usepackage{longtable,booktabs,array}
\usepackage{calc} % for calculating minipage widths
% Correct order of tables after \paragraph or \subparagraph
\usepackage{etoolbox}
\makeatletter
\patchcmd\longtable{\par}{\if@noskipsec\mbox{}\fi\par}{}{}
\makeatother
% Allow footnotes in longtable head/foot
\IfFileExists{footnotehyper.sty}{\usepackage{footnotehyper}}{\usepackage{footnote}}
\makesavenoteenv{longtable}
\usepackage{graphicx}
\makeatletter
\newsavebox\pandoc@box
\newcommand*\pandocbounded[1]{% scales image to fit in text height/width
  \sbox\pandoc@box{#1}%
  \Gscale@div\@tempa{\textheight}{\dimexpr\ht\pandoc@box+\dp\pandoc@box\relax}%
  \Gscale@div\@tempb{\linewidth}{\wd\pandoc@box}%
  \ifdim\@tempb\p@<\@tempa\p@\let\@tempa\@tempb\fi% select the smaller of both
  \ifdim\@tempa\p@<\p@\scalebox{\@tempa}{\usebox\pandoc@box}%
  \else\usebox{\pandoc@box}%
  \fi%
}
% Set default figure placement to htbp
\def\fps@figure{htbp}
\makeatother





\setlength{\emergencystretch}{3em} % prevent overfull lines

\providecommand{\tightlist}{%
  \setlength{\itemsep}{0pt}\setlength{\parskip}{0pt}}



 


% ============================================================
% bidi-auto.tex  —  XeLaTeX template for plain-text mixed
% Latin / Hebrew input with ZERO markup in the source.
%
% The Unicode Bidirectional Algorithm (UBA) handles direction
% automatically. XeLaTeX + bidi respect it at the paragraph
% level. Font fallback via fontspec covers missing glyphs.
%
% Compile:
%   pandoc input.md -o output.pdf \
%     --pdf-engine=xelatex \
%     --template=bidi-auto.tex
%
% Or standalone:
%   xelatex bidi-auto.tex
% ============================================================

% \documentclass[12pt, a4paper]{article}

\usepackage{iftex}
\RequireXeTeX

% ------------------------------------------------------------
% SINGLE FONT covering Latin + Hebrew — the key to zero markup.
% Taamey Frank CLM or SBL BibLit both span the full range.
% fontspec will fall back to the FallbackFonts list for any
% glyph the main font lacks; order matters.
% ------------------------------------------------------------
% \usepackage{fontspec}

% \setmainfont{FreeSans}[
%   Ligatures    = TeX,
%   FallbackFonts = {SBL BibLit, SBL Hebrew, FreeSans, DejaVu Sans}
% ]

% That's it for fonts. No \newfontfamily, no \texthebrew{}.

% ------------------------------------------------------------
% BIDI — loads the Unicode Bidirectional Algorithm support.
% With bidi loaded, XeLaTeX reads the Unicode script properties
% of each character and sets run direction automatically.
% No \LR or \RL commands needed in the source.
% ------------------------------------------------------------
\usepackage{fontspec}
\defaultfontfeatures{Renderer=HarfBuzz}
\setmainfont{Source Sans 3}[
  Ligatures    = TeX,
  FallbackFonts = {Alef, Arial Hebrew, Lucida Grande, FreeSans}
]
% \usepackage{fontspec}
% \defaultfontfeatures{Renderer=HarfBuzz}

% \setmainfont{Avenir Next}

% \newfontfamily\hebrewfont{Arial Hebrew}[
%   Renderer = HarfBuzz,
%   Script   = Hebrew
% ]

% \usepackage{fontspec}
% \defaultfontfeatures{Renderer=HarfBuzz}

% \setmainfont{Avenir Next}

\newfontfamily\hebrewfont{Alef}[
  Renderer = HarfBuzz,
  Script   = Hebrew
]

\usepackage[Hebrew]{ucharclasses}
% \setTransitionsFor{Hebrew}{\hebrewfont}{\normalfont}
% \usepackage[Hebrew, IPAExtensions]{ucharclasses}

\newfontfamily\ipafont{Source Sans 3}[Renderer=HarfBuzz]

\setTransitionsFor{Hebrew}{\hebrewfont}{\normalfont}
% \setTransitionsFor{IPAExtensions}{\ipafont}{\normalfont}
% \usepackage{ucharclasses}
% \setTransitionsForHebrew{\hebrewfont}{\normalfont}
\usepackage{bidi}
% \renewcommand{\,}{\thinspace}
% \usepackage{etoolbox}
% \AtBeginEnvironment{footnote}{\hbadness=10000\hfuzz=10pt}
% \let\oldfootnote\footnote
% \renewcommand{\footnote}[1]{\oldfootnote{\LTR #1}}
% \usepackage{etoolbox}
% \patchcmd{\footnote}{\ignorespaces}{\leavevmode\LTR\ignorespaces}{}{}
% \let\bidi@saved@everypar\everypar
% \usepackage{endfloat}
% \AtEndDocument{\let\LTR\relax\let\RTL\relax}
% \AtBeginDocument{
%   \let\LTR\relax
%   \let\RTL\relax
% }
\usepackage[hyperfootnotes=false]{hyperref}
% ============================================================
% header.tex  —  include-in-header snippet
% Replaces $var$ Pandoc syntax with \ifdefined\doc... guards.
% Vars injected via vars.lua Lua filter from YAML front matter:
%   \doctitle      \docsubtitle   \docauthor   \docdate
%   \docmeta       \docfoot       \doclogoleft \doclogoright
%
% No fancyhdr. Header fires via \AtBeginDocument,
% footer via \AtEndDocument — bidi-safe.
% ============================================================

\usepackage{graphicx}
\usepackage{tabularx}
\usepackage{xcolor}
\usepackage{array}
\usepackage{tikz}

% ---------- header font (adjust to taste) ----------
\newfontfamily\headerfont{Gill Sans Light}

% ---------- footer colors ----------
\definecolor{footerbg}{HTML}{000000}
\definecolor{footerfg}{HTML}{FFFFFF}

% ---------- column types ----------
\newcolumntype{L}{>{\raggedright\arraybackslash}m{10mm}}
\newcolumntype{C}{>{\centering\arraybackslash}X}
\newcolumntype{R}{>{\raggedleft\arraybackslash}m{45mm}}

% vertical separator rule
\newcommand{\vsep}{\color{black}\vrule width 1pt}

% ============================================================
% CUSTOM HEADER
% ============================================================
\newcommand{\CustomHeader}{%
\LTR{%
\noindent
\begin{tabularx}{\textwidth}{L !{\vsep} C !{\vsep} R}

% LEFT LOGO
\ifdefined\doclogoleft
  \includegraphics[width=10mm,height=10mm,keepaspectratio]{\doclogoleft}%
\fi
&
% CENTER TITLE BLOCK
\ifdefined\doctitle
  {\Large\bfseries\doctitle}%
\fi
\ifdefined\docsubtitle
  \\\normalsize\docsubtitle
\fi
&
% RIGHT LOGO
\ifdefined\doclogoright
  \includegraphics[height=10mm,keepaspectratio]{\doclogoright}%
\fi

\end{tabularx}

\hrule
\vspace{6pt}

% ---------- META AREA ----------
{\small
\ifdefined\docauthor
  \textbf{Author:} \docauthor\\
\fi
\ifdefined\docdate
  \textbf{Date:} \docdate
\fi
\ifdefined\docmeta
  \\\docmeta
\fi
}

\vspace{12pt}
}% end \LTR
}% end \CustomHeader

% ============================================================
% CUSTOM FOOTER  (overlay, anchored to page bottom)
% ============================================================
\newcommand{\CustomFooter}{%
\begin{tikzpicture}[remember picture, overlay]
\node[
  anchor=south,
  draw,
  fill=footerbg,
  text=footerfg,
  inner sep=10pt,
  rounded corners=2pt,
  text width=\dimexpr\textwidth-20pt\relax
]
at ([yshift=5mm]current page.south)
{%
  \begin{center}
  {\small
  \ifdefined\docfoot
    {\headerfont\bfseries\docfoot}%
  \fi
  }
  \end{center}
};
\end{tikzpicture}%
}% end \CustomFooter

% ============================================================
% FIRE header at doc start, footer at doc end
% ============================================================
% \AtBeginDocument{\CustomHeader}
\AtEndDocument{\CustomFooter}
\renewcommand{\maketitle}{}
\KOMAoption{captions}{tableheading}
\newcommand{\doctitle}{12293.wolf gossip satz für satz}
\newcommand{\docdate}{2026-03-22}
\newcommand{\docmeta}{class: [16412] dynamiken postdeutscher gegenwartslyrik :: tutor: chiara liso :: term: SS22 :: << back to paper <<}
\newcommand{\docfoot}{the anarkkiv lankalevikki worksheet collection >> https://ada-sub.dh-index.org/blog}
\newcommand{\doclogoleft}{../fadenlogo-fwb.png}
\newcommand{\doclogoright}{../dcl-head-niemand.png}
\newcommand{\docauthor}{stschwarz}
\makeatletter
\@ifpackageloaded{caption}{}{\usepackage{caption}}
\AtBeginDocument{%
\ifdefined\contentsname
  \renewcommand*\contentsname{Table of contents}
\else
  \newcommand\contentsname{Table of contents}
\fi
\ifdefined\listfigurename
  \renewcommand*\listfigurename{List of Figures}
\else
  \newcommand\listfigurename{List of Figures}
\fi
\ifdefined\listtablename
  \renewcommand*\listtablename{List of Tables}
\else
  \newcommand\listtablename{List of Tables}
\fi
\ifdefined\figurename
  \renewcommand*\figurename{Figure}
\else
  \newcommand\figurename{Figure}
\fi
\ifdefined\tablename
  \renewcommand*\tablename{Table}
\else
  \newcommand\tablename{Table}
\fi
}
\@ifpackageloaded{float}{}{\usepackage{float}}
\floatstyle{ruled}
\@ifundefined{c@chapter}{\newfloat{codelisting}{h}{lop}}{\newfloat{codelisting}{h}{lop}[chapter]}
\floatname{codelisting}{Listing}
\newcommand*\listoflistings{\listof{codelisting}{List of Listings}}
\makeatother
\makeatletter
\makeatother
\makeatletter
\@ifpackageloaded{caption}{}{\usepackage{caption}}
\@ifpackageloaded{subcaption}{}{\usepackage{subcaption}}
\makeatother
\usepackage{bookmark}
\IfFileExists{xurl.sty}{\usepackage{xurl}}{} % add URL line breaks if available
\urlstyle{same}
\hypersetup{
  pdfauthor={stschwarz},
  colorlinks=true,
  linkcolor={blue},
  filecolor={Maroon},
  citecolor={Blue},
  urlcolor={Blue},
  pdfcreator={LaTeX via pandoc}}


\title{12293.wolf gossip satz für satz}
\author{stschwarz}
\date{2026-03-22}
\begin{document}
\maketitle

\CustomHeader


\section{password parsley}\label{password-parsley}

der \textbf{login} auf einer der geschützten seiten einer mit unserer
universität verbundenen (institution) läuft über ein verfahren, das als
\textbf{single sign on} bezeichnet wird. bei diesem werden der
nutzername und das passwort (unsere zedat \textbf{credentials}) nicht
der von uns gewünschten seite übermittelt, sondern nur die zusage einer
instanz (die ich im folgenden näher beleuchten will), dasz unsere
credentials richtig sind, wir also der instanz (ich nenne sie hier
einmal: wächter) den richtigen nutzernamen und passwort übergeben haben,
um auf die ferne seite (ein korpusengine, ein bibliotheksportal,
generell eine ferne anwendung im netz, die den zugriff zb. nur den
\textbf{affiliates} bestimmter bildungseinrichtungen erlaubt / ich nenne
sie hier zb.: waffenlager, repositorium), zugreifen zu können. beide
seiten, wächter und repositorium, sind sich einig, dasz uns der zugriff
auf die ressourcen nur gestattet wird, wenn wir derjenige sind, der wir
zu sein vorgeben. zb. student/ der FUB.

\begin{itemize}
\tightlist
\item
  oder partisan.
\item
  oder mitglied einer geheimen verbindung.
\item
  oder: mitwisser von geheimnissen.
\end{itemize}

genauer: kenner des SHIBBOLETH. dieses beschäftigt hier also auch uljana
wolf.

\paragraph{1}\label{section}

wer gelegentlich ein biszchen in der literatur umherstreift, wird über
\textbf{codes \& switches} schon gestolpert sein, die nur denen
zugänglich sind, die dort irgendwie zuhause sind. die sich mit texten
auseinandersetzen, texte gegeneinander lesen -- auf ihre
intertextualität hin untersuchen, zb. / schnell macht man tropen aus,
die gemeinplätze geworden sind und mit denen jeder umgehen kann, der mit
dem kanon der jeweiligen (literatur) vertraut ist. einen solchen (switch
character) scheint hier, bei uljana, auch die petersilie zu haben. man
wird mit dem wort nicht weiter unbefangen hantieren können, wenn man
über die 20.000 haitianischen gastarbeiter
(cf.~wolf/GOSSIP/etymologischer gossip im gedicht)\footnote{Wolf, U.
  Etymologischer Gossip im Gedicht. In: Braun, M., \& Thill, H. (2018).
  Aus Mangel an Beweisen\,: deutsche Lyrik 2008-2018 / herausgegeben von
  Michael Braun und Hans Thill. Heidelberg: Wunderhorn.} gelesen hat. es
ist schön, dasz uljana hier die sache einmal aufklärt, um an diesem
beispiel das prinzip des \textbf{shibboleth} zu erklären.

das wort petersilie - parsley - {[}perejil{]} wird als zugang benötigt,
um gegenüber der gemeinschaft der natives NICHT als fremder zu
erscheinen = dem massaker zum opfer zu fallen. wer es {[}pɛʟɛχɪʟ{]}
ausspricht, und nicht {[}pɛʀɛχɪʟ{]}, wie die einheimischen, wird
ermordet. es ist also gut, die aussprache zu wissen. oder zu
beherrschen.

ähnliches, weniger drastisches, kann menschen widerfahren, die um nur
prominente beispiele zu nennen, {[}derrida{]}, {[}bourdieu{]},
{[}accessoir{]} oder eben, um zurückzukehren {[}shibboleth{]} (heute!
dazu siehe Ri12.6 unten.) falsch aussprechen. (cf.~dembeck,
pg.212).\footnote{Dembeck, T., Parr, R., \& Küpper, T. (2017). Literatur
  und Mehrsprachigkeit\,: ein Handbuch / Till Dembeck/Rolf Parr
  (Hrsg.)\,; unter Mitarbeit von Thomas Küpper. Tübingen: Narr Francke
  Attempto.}

\paragraph{2}\label{section-1}

S: שִׂבֹלֶת \textgreater{}
\href{https://www.pealim.com/dict/3270-shibolet/}{die übersetzung}
allein bringt niemanden weiter, es geht um die kleinen punkte\ldots, die
unscheinbaren merkzeichen, die das wort
\href{https://www.sefaria.org/Judges.12.6?lang=bi&with=all&lang2=en}{Ri,
12.6} an sich trägt und es zu dem machen, was es ist. nicht torheit,
sondern ein strom. das würde zu weit gehen, aber ob das Sibboleth doch
ein TOR ist in der deutsch-assoziierenden polyphonischen
auslegung\ldots, das hat sicher auch fr. wolf beschäftigt\ldots{} die
stelle in Ri12.6 jdfs. gibt auch dieser interpretation einen raum. die
richtige aussprache (hier nicht durch haitianische gastarbeiter, sondern
die dem gilead fremden ephraimiten) öffnet tore, verhindert, ermordet zu
werden.

\paragraph{3}\label{section-2}

in der literatur(wissenschaft) kann es essentiell werden, namen,
termini\ldots{} nicht nur richtig einzusetzen, sondern sie auch richtig
auszusprechen. dasz dabei hegemoniale\ldots{} (beweggründe? aspekte?
belange? sic!) zum tragen kommen, sei hier nicht nur erwähnt. die
\textbf{domaininhaber} definieren die korrektheit der anwendung. wenn
die dömäne (die dominikanische republik? der gilead?) ist, haben
diejenigen die epistemische? begriffshoheit, die dort \textbf{schon
immer ansässig sind}. in der literatur und ihrer wissenschaft die
worteigner und wissenschaftler der jeweiligen domäne. man setzt sich
schnell in irgendwelche nesseln, wenn man das shibboleth nicht kennt.
man fällt auf, aber man geht unter.

also kucken wir: was heiszt :hegemonial: was heiszt :epistemisch:. wir
kucken überall, wo nesseln warten könnten\ldots{} hegemonial im gossip,
(cf.~DIE WESTSÄULENLIEBHABEREI DER ÜBERSETZUNG. Ilse Aichingers
`schwache Architekur')\footnote{Prammer, T., Vescoli, C., \& Erb, E.
  (2021). Was für Sätze\,: Zu Ilse Aichinger / Theresia Prammer,
  Christine Vescoli (Hg.)\,; Elke Erb, Uljana Wolf, Monika Rinck, Ruth
  Klüger, Marlene Streeruwitz {[}und weitere{]}. (1. Auflage). Wien:
  Edition Korrespondenzen.}

\begin{quote}
Antihegemoniale, antikoloniale Übersetzung ist nicht der Macht treu, dem
Besten, dem Brutalen, sondern besteht auf diesen schwächeren, den `nicht
ausreichenden' Möglichkeiten. Eine Übersetzung muss untreu sein, um das
bloßzustellen, was das Beste `verbirgt'.
\end{quote}

zu :epistemisch: sagt fr. wolf nichts, aber (gramling, 2016)\footnote{Gramling,
  D. (2016). Zur Mehrsprachigkeitsforschung in der interkulturellen
  Literaturwissenschaft: Wende, Romanze, Rückkehr? Zeitschrift für
  interkulturelle Germanistik, 7(1), 133--150.
  \url{https://doi.org/10.14361/zig-2016-0109}}

\begin{quote}
In ihrer bescheidenen Weise vermochte es die Einsprachigkeit, zur Basis
für einen enormen Überbau ästhetischer und epistemischer Paradigmen zu
avancieren, denen man sich nur schwer entziehen kann -- einschließlich
des modernen Buchs, wie wir es kennen, und eines Kanons der
Weltliteratur, der das einsprachige (übersetzte oder übersetzbare) Buch
als Grundeinheit verwendet. Auch in den bildstürmerischsten Bereichen
der kulturellen Produktion -- sei es im Surrealismus, in der
Dekonstruktion, in der Kybernetik oder im Anarchismus -- spielt die
Einsprachigkeit eine zentrale Rolle. Sie bestimmt, was verlässlich
kommuniziert, verkehrsfähig übersetzt, politisch operationalisiert,
international verteilt und auch privat im Gedächtnis behalten werden
kann.
\end{quote}

können wir jetzt diese wörter benutzen? kennen wir nun ihre bedeutung?
versuchen Sie es mit :kongruenz:, :kontingenz:, :kohärenz:. nahe
beieinanderliegende phoneme? wir kucken nochmal: (Meibauer, Einführung
in die germanistische Linguistik, pg. 82)\footnote{Meibauer, J., Demske,
  U., Geilfuß-Wolfgang, J., Pafel, J., Ramers, K. H., Rothweiler, M., \&
  Steinbach, M. (2015). Einführung in die germanistische Linguistik (3.,
  überarbeitete und aktualisierte Auflage).
  \url{https://doi.org/10.1007/978-3-476-05424-1}}

\begin{quote}
Die Phoneme in (18b) stehen in Opposition bzw. sie kontrastieren.
Während Phone (Laute) in der Notation in eckigen Klammern stehen, werden
Phoneme in Schrägstriche eingefasst.
\end{quote}

gemeint waren also die \textbf{\emph{:phone:}} {[}kongruenz{]},
{[}kontingenz{]}, {[}kohärenz{]}. ein paar nesseln weniger\ldots{}

wir könnten das alles ohne weiteres fortschreiben; nur zu welchem
ausgang wollen wir damit gelangen? reicht der ansatz, um gewisztheit
über zumindest fr. wolfs engagement für ambiguitäten, ambivalenzen udgl.
{[}nicht die begriffe, sondern das bestreben UM{]} zu erlangen, wollte
man einer offenen literaturwissenschaft, entkanonisierten vielleicht
entkolonialisierten und\ldots, da sind wir: post{[}nationalen{]}{[}hier:
deutschen{]} literatur {[}et ici me manque encore une foi la verbe
finite allemand{]} u.u.: zum leben verhelfen. die gegend um fr. wolf
bietet jdfs. genug material, dies bestreben längst in taten umgesetzt zu
sehen.




\end{document}
